% !TeX program = pdflatex
% Strategia HR Basilisk Systems P.S.A. --- dostosowana do struktury prawnej spółki (umowa spółki 3.4-S, umowa akcjonariuszy 1.8, wzory umów).
% Źródło: strategia_hr.md (projekt do decyzji założycieli, listopad 2026). Odesłania pobierane z psa.aux i basilisk_sha.aux (xr-hyper).
\documentclass[11pt,a4paper]{report}
\usepackage[T1]{fontenc}
\usepackage[utf8]{inputenc}
\usepackage[english,polish]{babel}
\usepackage{xr-hyper}
\usepackage{basilisk_i18n}
\usepackage{lmodern}
\usepackage{microtype}
\usepackage{setspace}
\usepackage{parskip}
\usepackage{needspace}
\usepackage{tabularx}
\usepackage{float}
\externaldocument{psa}[psa.pdf]
\externaldocument{basilisk_sha}[basilisk_sha.pdf]

\BusinessPlanCompany{Basilisk Systems}
\BusinessPlanWatermarkText{Basilisk Systems}
\BusinessPlanWatermarkColor{black!6}
\BusinessPlanWatermarkFontSize{38}{46}
\BusinessPlanWatermarkAngle{38}
\BusinessPlanAuthorsText{Basilisk Systems --- Założyciele}
\BusinessPlanConfidentialText{Projekt do decyzji Założycieli --- poufne}
\BusinessPlanFooterText{Basilisk Systems --- strategia HR (projekt)}
\BusinessPlanVersionText{Wersja 1.0 --- projekt}

\setstretch{1.1}
\setlist[itemize]{topsep=0.2em,itemsep=0.12em,parsep=0.05em,leftmargin=1.8em}
\setlist[enumerate]{topsep=0.2em,itemsep=0.15em,parsep=0.05em,leftmargin=2.0em}
\setlength{\parindent}{0pt}
\setlength{\parskip}{0.5em}
\setlength{\emergencystretch}{4em}
\renewcommand{\arraystretch}{1.2}
\renewcommand{\thesection}{\arabic{section}}
\captionsetup{font=small,labelfont=bf}

\definecolor{cHR}{HTML}{1F4E79}
\newtcolorbox{Prawo}{colback=cHR!5, colframe=cHR!70, boxrule=0.6pt, arc=2pt, left=6pt, right=6pt, top=4pt, bottom=4pt, fonttitle=\small\bfseries, title=Podstawa prawna i dokumenty}
% style diagramów --- jak w podręczniku podejmowania decyzji
\definecolor{cWZ}{HTML}{1F4E79}
\definecolor{cRD}{HTML}{2E7D32}
\definecolor{cDW}{HTML}{8D6E00}
\definecolor{cX}{HTML}{8E2430}
\definecolor{cDHR}{HTML}{6A1B9A}
\tikzset{
  >={Stealth[length=2.2mm]},
  every node/.style={font=\small, align=center},
  box/.style={draw, rounded corners=2pt, minimum height=8mm, inner sep=4pt, text width=#1},
  box/.default=3.6cm,
  wz/.style={box=#1, fill=cWZ!12, draw=cWZ},
  wz/.default=3.6cm,
  rd/.style={box=#1, fill=cRD!12, draw=cRD},
  rd/.default=3.6cm,
  dw/.style={box=#1, fill=cDW!14, draw=cDW},
  dw/.default=3.6cm,
  xx/.style={box=#1, fill=cX!10, draw=cX},
  xx/.default=3.6cm,
  q/.style={draw, diamond, aspect=2.4, inner sep=1.5pt, fill=black!4, text width=#1},
  q/.default=3.4cm,
  lbl/.style={font=\scriptsize, fill=white, inner sep=1.5pt},
  line/.style={->, thick},
}
\newtcolorbox{DlaHR}{colback=cDHR!5, colframe=cDHR!75, boxrule=0.7pt, arc=2pt, left=6pt, right=6pt, top=4pt, bottom=4pt, fonttitle=\small\bfseries, title=Dla zespołu HR}
\newtcolorbox{Decyzja}{colback=black!3, colframe=black!45, boxrule=0.5pt, arc=2pt, left=6pt, right=6pt, top=4pt, bottom=4pt}

\newcommand{\field}[1]{\textbf{[#1]}}
\newcommand{\US}[1]{US §~#1}
\newcommand{\UA}[1]{UA §~#1}

\hypersetup{
  pdftitle={Basilisk Systems P.S.A. --- strategia HR},
  pdfauthor={Basilisk Systems}
}

\begin{document}

\BusinessPlanTitlePage
  {Basilisk Systems}
  {Wersja 1.0 --- projekt}
  {Strategia HR}
  {Organizacja pracy, rekrutacja i rozwój zespołu}
  {Dostosowana do umowy spółki, umowy akcjonariuszy i wzorów umów Basilisk Systems P.S.A.}
  {1.0}

\tableofcontents
\newpage

% ================================================================
\section{Cel strategii}

Basilisk powinien samodzielnie definiować wymagania, oceniać cudzą pracę, wydawać produkt i utrzymywać go u klienta. Sprawdzianem tej zdolności będzie \textbf{drugie wdrożenie}: czy inżynier zatrudniony spoza grona Założycieli potrafi je poprowadzić na podstawie kodu, dokumentacji i doświadczenia zgromadzonych w Spółce?

Dokument przenosi projekt strategii HR na strukturę prawną Spółki. Przy każdej roli wskazuje rodzaj umowy, wzór z pakietu dokumentów Spółki, organ, który zatwierdza wydatek, oraz sposób zabezpieczenia praw do wyników pracy. Skróty: „US” --- umowa spółki (wersja 3.4-S), „UA” --- umowa akcjonariuszy (wersja 1.8). Diagramy używają tych samych kolorów co podręcznik podejmowania decyzji, a fioletowe ramki „Dla zespołu HR” zawierają praktyczne wskazówki: listy dokumentów, terminy i sposób postępowania.

\begin{Decyzja}
\textbf{Wniosek główny: strategię można wdrożyć bez zmiany umowy spółki.} Umowa spółki nie reguluje struktury zespołu, stanowisk ani form zatrudnienia. Wystarczą decyzje Rady Dyrektorów (budżet, regulamin Rady, pełnomocnictwa) i --- w kilku sprawach --- uchwały akcjonariuszy przewidziane już w umowie spółki. Szczegóły i wyjątki: rozdział~\ref{sec:zmiany}.
\end{Decyzja}

% ================================================================
\section{Zasady wspólne dla wszystkich ról}
\label{sec:zasady}

\begin{enumerate}
  \item \textbf{Rola techniczna to nie umocowanie.} Nazwy ról (Opiekun Systemów, lider Ekipy, Ekspert) określają odpowiedzialność za wynik, a nie prawo do reprezentowania Spółki. Umowy podpisują dwaj dyrektorzy albo dyrektor z prokurentem (\US{\ref{art:reprezentacja} ust.~1}). Lider może ustalać priorytety pracy, ale wydatek zatwierdza się według budżetu i progów (\US{\ref{art:rada} ust.~7}).
  \item \textbf{O rodzaju umowy decydują rzeczywiste warunki pracy.} Praca pod kierownictwem Spółki, w wyznaczonym miejscu i czasie, jest zatrudnieniem na podstawie stosunku pracy, bez względu na nazwę umowy (art.~22 §~1\textsuperscript{1} KP). Ta sama ocena dotyczy regularnie nadzorowanych zmian testerów.
  \item \textbf{Prawa do wyników pracy przed pierwszym dniem pracy.} Pracownik --- z mocy prawa i umowy o pracę (art.~12 i 74 ust.~3 PrAut); każda inna osoba --- przez Umowę nr~2 w formie pisemnej, z kwartalnym potwierdzeniem wykazu utworów (\US{\ref{art:ip} ust.~1}).
  \item \textbf{Dostęp na czas i w zakresie zadania.} Dostępy nadaje i odbiera Odpowiedzialny według polityki bezpieczeństwa informacji; informacje Kontrolowane --- wyłącznie osobom z listy zatwierdzonej przez Radę. Osoba spoza Unii Europejskiej --- po ocenie wymogów kontroli eksportu i sprawdzeniu list sankcyjnych.
  \item \textbf{Poufność przed rozmową.} Kandydatom, partnerom i testerom informacje Poufne udostępnia się dopiero po zawarciu NDA albo umowy z klauzulą poufności (\UA{\ref{sha:poufnosc}}).
  \item \textbf{Jawność wynagrodzenia.} Kandydat na pracownika otrzymuje informację o wynagrodzeniu albo jego przedziale najpóźniej przed rozmową kwalifikacyjną, jeżeli nie podano jej w ogłoszeniu (art.~18\textsuperscript{3ca} KP).
  \item \textbf{Akcje dla zespołu tylko przez program motywacyjny.} Pracownikom, współpracownikom i Ekspertom akcje przyznaje się w programie motywacyjnym (seria~P, \US{\ref{art:seria-p}}), a nie przez obietnice w umowach o pracę. Seria~F jest zarezerwowana dla osób o roli porównywalnej z rolą Założyciela (\US{\ref{art:seria-f} ust.~1}).
\end{enumerate}

% ================================================================
\section{Siedem sposobów organizacji pracy}

Przy pięciu Założycielach jedna osoba może pełnić kilka funkcji. Nazwy pomagają ustalić, kto odpowiada za wynik i do kogo zwrócić się o pomoc. Osobny zespół warto utworzyć wtedy, gdy dana praca regularnie zajmuje czas kilku osób.

\begin{center}
\small
\begin{tabularx}{\linewidth}{@{}>{\raggedright\arraybackslash}p{3.0cm} >{\raggedright\arraybackslash}X >{\raggedright\arraybackslash}p{3.6cm} >{\raggedright\arraybackslash}p{3.0cm}@{}}
\toprule
\textbf{Rola} & \textbf{Podstawa współpracy} & \textbf{Prawa do wyników} & \textbf{Kto zatwierdza koszt} \\
\midrule
Opiekunowie Systemów (Core Keepers) & umowa o pracę z dodatkiem funkcyjnym i dyżurem (art.~151\textsuperscript{5} KP); Założyciel --- w ramach roli z \UA{\ref{sha:zaangazowanie}} & z mocy prawa / Umowa nr~2 & budżet: dyrektor wykonawczy \\
Kuźnia Inżynieryjna (Forge) & umowa o pracę & z mocy prawa & budżet \\
Ekipy Produktowe (Product Crews) & umowa o pracę; płatne zadanie rekrutacyjne --- umowa o dzieło albo zlecenia & z mocy prawa; zadanie --- licencja albo Umowa nr~2 & budżet \\
Eksperci Basiliska (Fellows) & aneks do umowy o pracę (mandat, premia za rezultat); ekspert zewnętrzny --- umowa B2B & z mocy prawa / Umowa nr~2 & budżet; poza budżetem ponad 250\,000 zł --- Rada \\
Partnerzy Specjalistyczni (Outrigger) & umowa o dzieło albo o świadczenie usług (B2B), cena za rezultat & Umowa nr~2 + licencja na elementy wcześniejsze & do 250\,000 zł poza budżetem --- dyrektor wyk.; ponad --- Rada; ponad 500\,000 zł --- WZ \\
Ekipa od Zadań Otwartych (Hackathon Force) & pracownicy w ramach obowiązków; osoby zewnętrzne --- umowa zlecenia & z mocy prawa / Umowa nr~2 & limit z karty mandatu \\
Sieć Terenowa (Field Mesh) & umowa zlecenia na zadanie (karta testu); przy regularnych, nadzorowanych zmianach --- umowa o pracę & licencja albo przeniesienie praw do danych; RODO & budżet testów \\
\bottomrule
\end{tabularx}
\end{center}

\begin{figure}[H]
\centering
\begin{tikzpicture}[node distance=10mm and 8mm]
  \node[rd=3.4cm] (crew) {\textbf{Ekipy Produktowe}\\produkt od prototypu do utrzymania};
  \node[wz=3.0cm, above left=12mm and 2mm of crew] (forge) {\textbf{Kuźnia}\\narzędzia, symulatory, stanowiska};
  \node[wz=3.0cm, above right=12mm and 2mm of crew] (fell) {\textbf{Eksperci}\\grant na trudny problem};
  \node[dw=3.0cm, left=16mm of crew] (keep) {\textbf{Opiekunowie Systemów}\\wydania, dostępy, odtworzenie};
  \node[dw=3.0cm, right=16mm of crew] (hack) {\textbf{Zadania Otwarte}\\tydzień na hipotezę};
  \node[xx=3.0cm, below left=12mm and 2mm of crew] (part) {\textbf{Partnerzy}\\rezultat z zewnątrz (B2B)};
  \node[xx=3.0cm, below right=12mm and 2mm of crew] (field) {\textbf{Sieć Terenowa}\\płatne próby (zlecenie)};
  \draw[line] (forge) -- node[lbl, left, pos=0.4]{narzędzie} (crew);
  \draw[line] (fell) -- node[lbl, right, pos=0.4]{metoda} (crew);
  \draw[line, <->] (keep) -- node[lbl, above]{wydanie} (crew);
  \draw[line] (hack) -- node[lbl, above]{prototyp} (crew);
  \draw[line] (part) -- node[lbl, left, pos=0.6]{odbiór} (crew);
  \draw[line] (field) -- node[lbl, right, pos=0.6]{wyniki prób} (crew);
\end{tikzpicture}
\caption{Siedem ról wokół Ekip Produktowych. Kolory: \textcolor{cRD}{\textbf{zielony}} --- rdzeń produktu; \textcolor{cWZ}{\textbf{niebieski}} --- wsparcie techniczne (etat); \textcolor{cDW}{\textbf{złoty}} --- role czasowe i dyżury; \textcolor{cX}{\textbf{czerwony}} --- osoby spoza zatrudnienia (B2B, zlecenie).}
\end{figure}

\begin{DlaHR}
\textbf{Teczka dokumentów według formy współpracy.}
\begin{itemize}
  \item \textbf{Pracownik} (Opiekun, Kuźnia, Ekipa, Ekspert): umowa o pracę podpisana przed dopuszczeniem do pracy; skierowanie i orzeczenie z badań wstępnych; szkolenie BHP wstępne przed rozpoczęciem pracy; zgłoszenie do ZUS w ciągu 7 dni; informacja o warunkach zatrudnienia (art.~29 §~3 KP) w ciągu 7 dni; potwierdzenie zapoznania z polityką bezpieczeństwa, regulaminem antymobbingowym i informacją o monitoringu; wykaz utworów sprzed zatrudnienia (Załącznik A2); dane do PPK. Dokumenty trafiają do akt osobowych (części A--E).
  \item \textbf{Konsultant albo Partner (B2B):} umowa (wzór~B), \textbf{Umowa nr~2 podpisana odręcznie albo kwalifikowanym podpisem} (bez niej Spółka nie ma praw do wyników), NDA, karta zlecenia, oświadczenia z Załącznika B3. Spółka nie prowadzi akt osobowych ani nie zgłasza do ZUS.
  \item \textbf{Tester (zlecenie):} umowa ramowa, karta testu dla każdego zadania, ewidencja godzin, przeniesienie praw do danych, klauzula RODO; zgłoszenie do ZUS według statusu (np. student do 26.~roku życia --- bez składek).
\end{itemize}
\end{DlaHR}

\subsection*{3.1. Opiekunowie Systemów (Core Keepers)}
\textbf{Kto:} osoby z zespołu, które wykazały się rzetelnością przy wydaniach i usuwaniu awarii; droga otwarta także dla nowych pracowników. \textbf{Samodzielność:} prowadzą wydania, zarządzają dostępami i odtwarzają usługę według uzgodnionych procedur. \textbf{Wynagrodzenie:} dodatek za rolę i uzgodnione warunki dyżuru. \textbf{Ocena:} dyrektor odpowiedzialny za produkt sprawdza, czy zastępca umie przywrócić poprzednią wersję; otrzymuje zapis próby.
\begin{Prawo}
Dyżur pracownika (gotowość poza godzinami pracy) wymaga zachowania prawa do odpoczynku; za dyżur poza domem przysługuje czas wolny albo wynagrodzenie (art.~151\textsuperscript{5} KP). Warunki dyżuru wpisuje się do umowy o pracę albo regulaminu. „Zarządzanie dostępami” oznacza wykonywanie polityki bezpieczeństwa informacji z upoważnienia Odpowiedzialnego --- rejestr dostępów i decyzje co do informacji Kontrolowanych pozostają przy Odpowiedzialnym i Radzie. Jeżeli Opiekunem jest Założyciel, jego rolę i minimalny zakres zaangażowania warto wpisać do tabeli w \UA{\ref{sha:zaangazowanie} ust.~1}.
\end{Prawo}

\subsection*{3.2. Kuźnia Inżynieryjna (Forge)}
\textbf{Kto:} specjaliści od integracji, automatyzacji i pomiarów, z otwartego naboru lub z zespołów produktowych. \textbf{Samodzielność:} wybierają metody budowy wspólnych narzędzi, symulatorów i stanowisk kalibracyjnych. \textbf{Wynagrodzenie:} stałe, z możliwością awansu technicznego. \textbf{Ocena:} inżynier z Ekipy Produktowej odtwarza wynik na podstawie instrukcji; Kuźnia utrzymuje narzędzie.
\begin{Prawo}
Umowa o pracę według wzoru~A z pakietu „Umowy z konsultantami”, z obowiązkami twórczymi w Załączniku A1. Narzędzia, symulatory i dokumentacja są utworami pracowniczymi --- prawa nabywa Spółka (art.~12 i 74 ust.~3 PrAut). Stanowiska kalibracyjne i sprzęt kupuje Spółka według progów (\US{\ref{art:rada} ust.~7}); wynalazki pracownicze --- art.~11 ust.~3 i 22 PWP.
\end{Prawo}

\subsection*{3.3. Ekipy Produktowe (Product Crews)}
\textbf{Kto:} otwarta rekrutacja przez ogłoszenia, uczelnie i społeczności techniczne; kandydat pokazuje własną pracę albo wykonuje krótkie płatne zadanie. \textbf{Samodzielność:} prowadzą produkt od prototypu przez instalację po utrzymanie; lider ustala kolejność prac w granicach mandatu projektu. \textbf{Wynagrodzenie:} etat z jawnym zakresem odpowiedzialności i kryteriami awansu. \textbf{Ocena:} osoba odpowiedzialna za wdrożenie potwierdza odbiór; Ekipa zachowuje historię decyzji, kontakt techniczny z klientem i obsługę zgłoszeń.
\begin{Prawo}
Umowa o pracę (wzór~A). \textbf{Płatne zadanie rekrutacyjne} zleca się na podstawie krótkiej umowy z ustaloną stawką, limitem czasu i zasadą wykorzystania materiałów: kandydat zachowuje prawa, a Spółka otrzymuje licencję do oceny; jeżeli Spółka chce wykorzystać wynik w produkcie --- przeniesienie praw w formie pisemnej. Przed zadaniem kandydat podpisuje NDA, a dane z zadania nie obejmują informacji Kontrolowanych. Kontakt techniczny z klientem nie obejmuje prawa do składania oświadczeń woli w imieniu Spółki.
\end{Prawo}

\subsection*{3.4. Eksperci Basiliska (Basilisk Fellows)}
\textbf{Kto:} kandydatura własna albo zgłoszenie przez zespół; decydują wyniki i propozycja rozwiązania konkretnego problemu. \textbf{Samodzielność:} chroniony czas, sprzęt i budżet na określony okres; ekspert wybiera metodę i współpracownika. \textbf{Wynagrodzenie:} za pracę oraz uzgodniona premia za rezultat; budżet eksperymentu pokrywa wydatki projektu. \textbf{Ocena:} sponsor ocenia wynik, a drugi inżynier odtwarza go samodzielnie; po przeglądzie --- zakończenie albo odnowienie grantu.
\begin{Prawo}
Grant ekspercki ma formę \textbf{aneksu do umowy o pracę} (albo do umowy B2B) z kartą mandatu (wzór~1). Premia za rezultat musi mieć obiektywne kryteria i termin --- premia regulaminowa jest roszczeniem pracownika, uznaniowa nie. Budżet eksperymentu mieści się w budżecie rocznym; wydatki poza budżetem powyżej 250\,000 zł --- uchwała Rady. Wybór współpracownika zewnętrznego --- według zasad dla Partnerów (pkt~3.5). Ekspertowi można dodatkowo przyznać opcje w programie motywacyjnym.
\end{Prawo}

\subsection*{3.5. Partnerzy Specjalistyczni (Outrigger Partners)}
\textbf{Kto:} niezależne firmy albo specjaliści dobrani na podstawie próbki prac, oferty i zdolności wykonania zlecenia. \textbf{Samodzielność:} organizują wykonanie uzgodnionego rezultatu własnymi metodami. \textbf{Wynagrodzenie:} cena za rezultat albo odebrane etapy. \textbf{Ocena:} imiennie wyznaczony inżynier Basiliska odbiera dostawę, dokumentację i uzgodnione prawa do używania i serwisowania rozwiązania.
\begin{Prawo}
Umowa o świadczenie usług doradczych (wzór~B) albo umowa o dzieło z kartą zlecenia (wzór~2), \textbf{Umowa nr~2} oraz licencja na elementy wcześniejsze partnera z prawem do sublicencji klientom --- bez niej Spółka nie może licencjonować produktu (art.~67 ust.~3 PrAut). Partner nie wprowadza komponentów na licencjach copyleft ani modeli i danych z ograniczeniem zastosowań wojskowych bez zgody Spółki; dostarcza wykaz komponentów z krajem pochodzenia i klasyfikacją eksportową. Inżynier odbierający \textbf{potwierdza odbiór merytoryczny}, ale umowę i protokół odbioru skutkujący zapłatą podpisują osoby uprawnione do reprezentacji albo pełnomocnik Rady. Umowa z podmiotem powiązanym z akcjonariuszem posiadającym co najmniej 5\% akcji albo z dyrektorem --- uchwała Rady, a ponad 100\,000 zł rocznie także Walnego Zgromadzenia (\US{\ref{art:wz-zastrzezone} ust.~1 lit.~k}).
\end{Prawo}

\subsection*{3.6. Ekipa od Zadań Otwartych (Hackathon Force)}
\textbf{Kto:} osoba albo duet z inicjatywą techniczną albo handlową; może powstać na tydzień z osób z różnych ekip. \textbf{Samodzielność:} w limicie czasu i wydatków wybiera sposób sprawdzenia jednej hipotezy. \textbf{Wynagrodzenie:} płatny czas eksperymentu; wartością jest także wiarygodny wynik uzasadniający zamknięcie pomysłu. \textbf{Ocena:} sponsor decyduje o dalszej inwestycji; przy kontynuacji Ekipa Produktowa dostaje dane, prototyp i opis ograniczeń.
\begin{Prawo}
Dla pracowników to zadanie w ramach obowiązków --- bez odrębnej umowy; karta mandatu (wzór~1) określa limit wydatków. Osoba zewnętrzna uczestniczy wyłącznie na podstawie umowy zlecenia z NDA i Umową nr~2. Prototypy i dane z eksperymentu zapisuje się w systemach Spółki; testy dronów w terenie --- z zachowaniem przepisów o bezzałogowych statkach powietrznych.
\end{Prawo}

\subsection*{3.7. Sieć Terenowa (Field Mesh)}
\textbf{Kto:} dobrowolny, płatny nabór testerów i osób zbierających dane, dobieranych według miejsca, urządzenia i warunków próby. \textbf{Samodzielność:} uczestnik wybiera dostępne zadanie i wykonuje je według karty testu. \textbf{Wynagrodzenie:} stawka znana przed przyjęciem zadania, także za poprawną próbę zakończoną awarią produktu. \textbf{Ocena:} koordynator sprawdza jakość próbki i uprawnienia do jej wykorzystania, rozlicza pracę i przekazuje wynik zespołowi produktu.
\begin{Prawo}
\begin{itemize}
  \item \textbf{Umowa:} umowa ramowa zlecenia z kartą testu (wzór~3) dla każdego zadania; minimalna stawka godzinowa i ewidencja godzin (ustawa o minimalnym wynagrodzeniu za pracę), składki według statusu testera. Przy regularnych zmianach, w wyznaczonym miejscu i czasie, pod nadzorem koordynatora --- ryzyko stosunku pracy (art.~22 §~1\textsuperscript{1} KP).
  \item \textbf{Prawa do danych:} tester przenosi na Spółkę prawa do zdjęć, nagrań i opisów albo udziela licencji wyłącznej, w formie pisemnej; jeżeli dane obejmują wizerunek lub dane osób trzecich --- podstawa przetwarzania i informacja według RODO, a w razie potrzeby zgoda na rozpowszechnianie wizerunku.
  \item \textbf{Bezpieczeństwo:} testerzy nie otrzymują informacji Kontrolowanych ani dostępu do kodu; urządzenia testowe są własnością Spółki, a testy z użyciem dronów prowadzi się w kategoriach i z uprawnieniami wymaganymi przez rozporządzenia (UE) 2019/945 i 2019/947. Spółka zapewnia ubezpieczenie OC dla prób w terenie.
  \item \textbf{Brak wzoru:} pakiet dokumentów Spółki nie zawiera jeszcze wzoru umowy zlecenia dla testerów --- należy go przygotować przed pierwszym naborem.
\end{itemize}
\end{Prawo}

% ================================================================
\section{Rekrutacja i rozwój}

Pierwsze ogłoszenie poza siecią Założycieli powinno opisywać rzeczywiste zadanie, np. doprowadzenie podsystemu od uruchomienia na stanowisku do instalacji u klienta. Należy podać zakres płacy (art.~18\textsuperscript{3ca} KP), warunki pracy, etapy wyboru i nazwisko przyszłego opiekuna. Nazwa stanowiska i treść ogłoszenia są neutralne pod względem płci (art.~18\textsuperscript{3ca} §~3 KP). Polecenia pracownicze i kandydatury z otwartego naboru ocenia się według tych samych kryteriów.

Na próbce pracy sprawdza się trzy rzeczy: czy kandydat potrafi znaleźć przyczynę problemu, wykazać poprawność wyniku i przekazać rozwiązanie drugiej osobie. Przed płatnym zadaniem uzgadnia się stawkę, limit czasu i zasady wykorzystania materiałów (pkt~3.3). Po zatrudnieniu nowa osoba najpierw odtwarza istniejące wydanie z mentorem, następnie prowadzi zmianę i uczestniczy w jej odbiorze.

\begin{figure}[H]
\centering
\begin{tikzpicture}[node distance=5mm and 5mm]
  \node[wz=3.0cm] (s1) {\textbf{1. Ogłoszenie}\\zadanie, przedział płacy, etapy, opiekun};
  \node[wz=3.0cm, right=of s1] (s2) {\textbf{2. Próbka pracy}\\NDA; płatne zadanie na umowę};
  \node[wz=3.0cm, right=of s2] (s3) {\textbf{3. Rozmowa}\\przyczyna, dowód, przekazanie};
  \node[rd=3.0cm, right=of s3] (s4) {\textbf{4. Decyzja}\\budżet; poza budżetem --- progi};
  \node[rd=3.0cm, below=12mm of s4] (s5) {\textbf{5. Umowa}\\wzór A albo B + prawa do wyników};
  \node[dw=3.0cm, left=of s5] (s6) {\textbf{6. Pierwszy dzień}\\BHP, badania, polityka, dostępy};
  \node[dw=3.0cm, left=of s6] (s7) {\textbf{7. Mentor}\\odtworzenie istniejącego wydania};
  \node[dw=3.0cm, left=of s7] (s8) {\textbf{8. Samodzielność}\\prowadzenie zmiany i odbiór};
  \draw[line] (s1) -- (s2); \draw[line] (s2) -- (s3); \draw[line] (s3) -- (s4);
  \draw[line] (s4) -- (s5); \draw[line] (s5) -- (s6); \draw[line] (s6) -- (s7); \draw[line] (s7) -- (s8);
\end{tikzpicture}
\caption{Od ogłoszenia do samodzielnej pracy: \textcolor{cWZ}{\textbf{rekrutacja}}, \textcolor{cRD}{\textbf{decyzja i umowa}}, \textcolor{cDW}{\textbf{wdrożenie}}.}
\end{figure}

Awans może prowadzić do zarządzania produktem albo do samodzielnej pracy eksperckiej. Pracownik, który rozwiązuje trudne problemy i uczy innych, może otrzymać \textbf{grant ekspercki}. Pracownik, który potrafi rozdzielić pracę, zaplanować odbiór i poprowadzić kontakt z klientem, może objąć kierownictwo wdrożenia. Każda osoba w zespole powinna znać kryteria i móc zgłosić kandydaturę.

\begin{Prawo}
\begin{itemize}
  \item \textbf{Kryteria awansu i premii} zapisuje się w dokumencie przyjętym przez Radę Dyrektorów. Regulamin wynagradzania i regulamin pracy są obowiązkowe od 50 pracowników (art.~77\textsuperscript{2} i 104 KP); wcześniej wystarczą zasady Rady, umowy o pracę i obwieszczenie o monitoringu.
  \item \textbf{Program motywacyjny:} opcje na akcje serii~P (pula 176\,471 akcji) według regulaminu programu; przydziały w granicach rocznego planu zatwierdzonego przez Walne Zgromadzenie.
  \item \textbf{Przetwarzanie danych kandydatów} --- w zakresie art.~22\textsuperscript{1} KP; dane z próbek pracy usuwa się po zakończeniu rekrutacji, chyba że kandydat zgodzi się na dłuższe przechowywanie.
\end{itemize}
\end{Prawo}

\begin{figure}[H]
\centering
\begin{tikzpicture}[node distance=8mm and 8mm]
  \node[rd=3.6cm] (inz) {\textbf{Inżynier w Ekipie Produktowej}};
  \node[q=3.2cm, below=of inz] (q) {Co robi najlepiej?};
  \node[wz=3.8cm, below left=10mm and 2mm of q] (e1) {\textbf{Ścieżka ekspercka}\\rozwiązuje trudne problemy i uczy innych $\to$ grant ekspercki};
  \node[dw=3.8cm, below right=10mm and 2mm of q] (m1) {\textbf{Ścieżka kierownicza}\\dzieli pracę, planuje odbiór, prowadzi klienta $\to$ kierownik wdrożenia};
  \node[wz=3.8cm, below=of e1] (e2) {\textbf{Ekspert Basiliska}\\kolejne granty po przeglądzie};
  \node[dw=3.8cm, below=of m1] (m2) {\textbf{Kierownik produktu}\\zespół i budżet w granicach mandatu};
  \node[box=8.4cm, fill=black!3, below=8mm of e2.south east, anchor=north west, xshift=-4.0cm] (op) {\textbf{Niezależnie od ścieżki:} opcje na akcje serii~P według regulaminu programu motywacyjnego};
  \draw[line] (inz) -- (q); \draw[line] (q) -| (e1); \draw[line] (q) -| (m1);
  \draw[line] (e1) -- (e2); \draw[line] (m1) -- (m2);
\end{tikzpicture}
\caption{Dwie równorzędne ścieżki awansu. Kryteria są jawne, a każda osoba może zgłosić kandydaturę.}
\end{figure}

\begin{DlaHR}
\begin{itemize}
  \item \textbf{Ogłoszenie:} przedział wynagrodzenia (albo przekazanie go przed rozmową), neutralna nazwa stanowiska (np. „Inżynier / Inżynierka systemów” albo „Inżynier (K/M)”), opis rzeczywistego zadania i etapów.
  \item \textbf{Dane kandydatów:} zbieramy tylko dane z art.~22\textsuperscript{1} KP (imię i nazwisko, data urodzenia, kontakt, wykształcenie, kwalifikacje, przebieg zatrudnienia). Nie prosimy o zdjęcie, PESEL ani informacje o rodzinie. Zgoda na udział w przyszłych rekrutacjach --- osobno. Po zakończeniu rekrutacji usuwamy CV i próbki pracy osób nieprzyjętych.
  \item \textbf{Płatne zadanie:} umowa (stawka, limit czasu, zasady wykorzystania materiałów) i NDA podpisane \emph{przed} przekazaniem zadania; zadanie nie zawiera informacji Kontrolowanych.
  \item \textbf{Ocena:} ten sam arkusz dla poleceń pracowniczych i kandydatów z otwartego naboru --- trzy kryteria: przyczyna problemu, dowód poprawności, przekazanie rozwiązania.
  \item \textbf{Pierwszy dzień:} zestaw powitalny z listą kontrolną (dokumenty, BHP, sprzęt, konta z 2FA, polityka, regulamin antymobbingowy, kontakt do mentora i Pełnomocnika).
\end{itemize}
\end{DlaHR}


% ================================================================
\section{Wybór umowy i granice odpowiedzialności}

\begin{itemize}
  \item \textbf{Umowa o pracę} --- stałe funkcje wykonywane pod kierownictwem Spółki, w wyznaczonym miejscu i czasie (wzór~A). Zakaz konkurencji wyłącznie w odrębnej umowie w formie pisemnej (wzór~C), po ustaniu zatrudnienia --- z odszkodowaniem co najmniej 25\% (art.~101\textsuperscript{2} KP). Kar umownych wobec pracowników nie stosuje się.
  \item \textbf{Niezależne podwykonawstwo} (wzór~B albo umowa o dzieło) --- gdy wykonawca organizuje pracę, a Basilisk umie z góry określić i zweryfikować rezultat. Przed zamówieniem wskazuje się osobę odpowiedzialną za odbiór. Gdy brakuje takiej kompetencji, najpierw powierza się jej rozwój Kuźni albo Ekspertowi.
  \item \textbf{Kontrakt menedżerski} --- przy rzeczywistym mandacie zarządczym (zespół, budżet, decyzje samodzielne, raportowanie, zastępstwo, zakończenie mandatu). Podstawę współpracy dobiera się do faktycznych warunków. Menedżer niebędący dyrektorem nie reprezentuje Spółki, chyba że otrzyma \textbf{pełnomocnictwo} albo prokurę (\US{\ref{art:reprezentacja} ust.~1--2}). Jeżeli menedżer jest dyrektorem, umowę z nim zawiera pełnomocnik powołany uchwałą akcjonariuszy (\US{\ref{art:reprezentacja} ust.~3}), a wynagrodzenie za funkcję dyrektora ustala Walne Zgromadzenie (\US{\ref{art:wz-zastrzezone} ust.~3}).
  \item \textbf{Role techniczne, dostęp i podpisy} rozpatruje się oddzielnie. Kierownik może ustalać priorytety pracy. Zatwierdzenie wydatku i podpisanie umowy następują według dokumentów Spółki: reprezentacja łączna (dwóch dyrektorów albo dyrektor z prokurentem), progi z \US{\ref{art:rada} ust.~7} i \US{\ref{art:wz-zastrzezone} ust.~1}, odrębny tryb dla umów z dyrektorami.
\end{itemize}

Przed rozpoczęciem pracy korzysta się z istniejących wzorów: umowy z konsultantami (wzory~A--C), Umowy nr~2, NDA i polityki bezpieczeństwa. Przy zleceniu zewnętrznym uzgadnia się prawa do wcześniejszych elementów wykonawcy i materiały potrzebne do utrzymania wyniku. Dostępy nadaje się na czas i w zakresie zadania, a przy zakończeniu współpracy stosuje się procedurę przy odejściu z polityki bezpieczeństwa (odebranie dostępów najpóźniej w ostatnim dniu, wydanie utworów w 14 dni).

\begin{figure}[H]
\centering
\begin{tikzpicture}[node distance=7mm and 10mm]
  \node[q=3.6cm] (q1) {Praca pod kierownictwem, w miejscu i czasie wyznaczonym przez Spółkę?};
  \node[rd=4.6cm, right=of q1] (a1) {\textbf{Umowa o pracę} (wzór A)\\zakaz konkurencji --- tylko wzór~C};
  \node[q=3.6cm, below=of q1] (q2) {Rezultat określony z góry i możliwy do odbioru?};
  \node[xx=4.6cm, right=of q2] (a2) {\textbf{B2B albo umowa o dzieło} (wzór B)\\+ Umowa nr~2 + karta zlecenia};
  \node[q=3.6cm, below=of q2] (q3) {Pojedyncze zadanie: test, zadanie rekrutacyjne, hackathon?};
  \node[dw=4.6cm, right=of q3] (a3) {\textbf{Umowa zlecenia}\\+ karta testu + przeniesienie praw};
  \node[wz=3.8cm, below=of q3] (a4) {Najpierw zdefiniuj rezultat (Kuźnia albo Ekspert)};
  \draw[line] (q1) -- node[lbl, above]{tak} (a1);
  \draw[line] (q2) -- node[lbl, above]{tak} (a2);
  \draw[line] (q3) -- node[lbl, above]{tak} (a3);
  \draw[line] (q1) -- node[lbl, left]{nie} (q2);
  \draw[line] (q2) -- node[lbl, left]{nie} (q3);
  \draw[line] (q3) -- node[lbl, left]{nie} (a4);
  \node[box=4.6cm, fill=black!3, below=10mm of a3] (dyr) {\textbf{Zawsze sprawdź:} druga strona jest dyrektorem $\to$ pełnomocnik z uchwały akcjonariuszy; akcjonariusz $\geq$5\% albo podmiot powiązany $\to$ Rada (+ WZ ponad 100\,000 zł rocznie)};
\end{tikzpicture}
\caption{Jaką umowę zawrzeć? O wyborze decydują rzeczywiste warunki pracy (art.~22 §~1\textsuperscript{1} KP).}
\end{figure}

\begin{DlaHR}
\textbf{Sygnały ostrzegawcze przy umowie B2B albo zleceniu} --- jeżeli występują razem, porozmawiaj z dyrektorem wykonawczym o umowie o pracę: stałe godziny wyznaczane przez Spółkę; przełożony wydaje bieżące polecenia co do sposobu pracy; osoba nie może powierzyć zadania nikomu innemu; rozliczenie za czas obecności, a nie za rezultat; udział w grafiku dyżurów na równi z pracownikami. Kontrola inspekcji pracy albo sąd mogą ustalić, że strony łączy stosunek pracy, z zaległymi składkami i uprawnieniami pracowniczymi.

\textbf{Obieg umowy:} projekt z właściwego wzoru $\to$ zatwierdzenie kosztu według progów $\to$ podpis dwóch dyrektorów albo dyrektora z prokurentem (przy umowie z dyrektorem --- pełnomocnik) $\to$ egzemplarz do akt i rejestru umów $\to$ zlecenie nadania dostępów Odpowiedzialnemu. Dział HR przygotowuje dokumenty, ale nie podpisuje umów w imieniu Spółki bez pełnomocnictwa.
\end{DlaHR}

\begin{figure}[H]
\centering
\begin{tikzpicture}[node distance=5mm and 5mm]
  \node[box=3.0cm, fill=black!8] (o1) {\textbf{Wypowiedzenie albo porozumienie}\\data ostatniego dnia};
  \node[rd=3.0cm, right=of o1] (o2) {\textbf{Ostatni dzień}\\odebranie dostępów, zwrot sprzętu, protokół};
  \node[rd=3.0cm, right=of o2] (o3) {\textbf{Do 14 dni}\\wydanie utworów, oświadczenie przy odejściu};
  \node[dw=3.0cm, right=of o3] (o4) {\textbf{Rozliczenie}\\świadectwo pracy, ZUS, dokumentacja};
  \node[xx=4.6cm, below=10mm of o2.south east, anchor=north] (zk) {\textbf{Zakaz konkurencji (wzór C)?}\\odszkodowanie w miesięcznych ratach przez okres zakazu};
  \draw[line] (o1) -- (o2); \draw[line] (o2) -- (o3); \draw[line] (o3) -- (o4);
  \draw[line, dashed] (o1) |- (zk);
\end{tikzpicture}
\caption{Zakończenie współpracy (polityka bezpieczeństwa informacji, procedura przy odejściu).}
\end{figure}

\begin{DlaHR}
\textbf{Kalendarz odejścia pracownika:} świadectwo pracy wydaje się w dniu ustania stosunku pracy (a gdy to niemożliwe --- w ciągu 7 dni), razem z informacją o okresie przechowywania dokumentacji pracowniczej (art.~94\textsuperscript{6} KP); wyrejestrowanie z ZUS w ciągu 7 dni; ekwiwalent za niewykorzystany urlop w ostatniej wypłacie; przy umowie o zakazie konkurencji --- harmonogram rat odszkodowania. Dodatkowo: rozmowa przy odejściu (anonimowo do badania atmosfery), oświadczenie według wzoru z polityki i potwierdzenie od Odpowiedzialnego, że dostępy zostały odebrane. Przy konsultantach i testerach: protokół wydania wyników, ostatnie potwierdzenie wykazu utworów (Umowa nr~2), odebranie dostępów.
\end{DlaHR}

\subsection*{Kto podpisuje i kto zatwierdza}

\begin{figure}[H]
\centering
\begin{tikzpicture}[x=0.0165cm, y=1cm]
  \draw[->] (0,0) -- (640,0) node[right, font=\scriptsize]{tys. zł};
  \foreach \v in {0,100,250,500} {\draw (\v,0.08) -- (\v,-0.08) node[below, font=\scriptsize]{\v};}
  \node[anchor=east, font=\scriptsize, text width=3.4cm, align=right] at (-8,1.2) {umowa / wydatek\\poza budżetem};
  \fill[cDW!25] (0,0.85) rectangle (250,1.55);  \node[font=\scriptsize] at (125,1.2) {dyrektor wykonawczy};
  \fill[cRD!25] (250,0.85) rectangle (500,1.55); \node[font=\scriptsize] at (375,1.2) {Rada Dyrektorów};
  \fill[cWZ!25] (500,0.85) rectangle (630,1.55); \node[font=\scriptsize] at (565,1.2) {WZ 75\%};
  \node[anchor=east, font=\scriptsize, text width=3.4cm, align=right] at (-8,2.3) {umowa z podmiotem\\powiązanym (rocznie)};
  \fill[cRD!25] (0,1.95) rectangle (100,2.65);  \node[font=\scriptsize] at (50,2.3) {Rada};
  \fill[cWZ!25] (100,1.95) rectangle (630,2.65); \node[font=\scriptsize] at (365,2.3) {Rada + WZ 75\%};
  \node[anchor=east, font=\scriptsize, text width=3.4cm, align=right] at (-8,3.4) {zatrudnienie, zlecenia\\w ramach budżetu};
  \fill[cDW!25] (0,3.05) rectangle (630,3.75); \node[font=\scriptsize] at (315,3.4) {dyrektor wykonawczy (delegacja Rady)};
  \foreach \v in {100,250,500} \draw[densely dotted, black!50] (\v,0) -- (\v,3.85);
\end{tikzpicture}
\caption{Kto zatwierdza koszt zatrudnienia albo zlecenia (\US{\ref{art:rada} ust.~7}, \US{\ref{art:wz-zastrzezone} ust.~1 lit.~h i~k}).}
\end{figure}

\begin{center}
\small
\begin{tabularx}{\linewidth}{@{}>{\raggedright\arraybackslash}X >{\raggedright\arraybackslash}p{4.4cm} >{\raggedright\arraybackslash}p{4.0cm}@{}}
\toprule
\textbf{Czynność} & \textbf{Zatwierdza} & \textbf{Podpisuje} \\
\midrule
Zatrudnienie albo zlecenie w ramach budżetu & dyrektor wykonawczy (delegacja Rady) & dwóch dyrektorów / dyrektor z prokurentem / pełnomocnik \\
Umowa poza budżetem do 250\,000 zł & dyrektor wykonawczy & jw. \\
Umowa poza budżetem 250\,000--500\,000 zł & Rada Dyrektorów & jw. \\
Umowa poza budżetem ponad 500\,000 zł & Walne Zgromadzenie 75\% & jw. \\
Umowa z akcjonariuszem $\geq$5\%, dyrektorem, osobą bliską albo podmiotem powiązanym & Rada; ponad 100\,000 zł rocznie --- także WZ 75\% & jw.; z dyrektorem --- pełnomocnik z uchwały akcjonariuszy \\
Wynagrodzenie dyrektora za pełnienie funkcji & Walne Zgromadzenie & --- \\
Przydział opcji (seria P) & Rada w granicach planu WZ; dyrektorowi --- WZ & jw.; z dyrektorem --- pełnomocnik \\
Dostęp do informacji Kontrolowanych & Rada Dyrektorów (lista imienna) & --- \\
\bottomrule
\end{tabularx}
\end{center}

% ================================================================
\section{Pierwsze płatne wdrożenie}

Przykładowy pilot obejmuje instalację oprogramowania i sprawdzenie jakości obrazu na urządzeniach klienta.

\begin{figure}[H]
\centering
\begin{tikzpicture}[node distance=5mm and 5mm]
  \node[rd=3.0cm] (p1) {\textbf{1. Obsada}\\Założyciel prowadzący, wydanie, zastępca, nowa osoba};
  \node[dw=3.0cm, right=of p1] (p2) {\textbf{2. Rozpoznanie}\\Zadania Otwarte (tydzień) + Sieć Terenowa};
  \node[wz=3.0cm, right=of p2] (p3) {\textbf{3. Rozwiązanie}\\Ekspert, Partner, odbiór w Kuźni};
  \node[rd=3.0cm, right=of p3] (p4) {\textbf{4. Przekazanie}\\druga instalacja przez nową osobę};
  \draw[line] (p1) -- (p2); \draw[line] (p2) -- (p3); \draw[line] (p3) -- (p4);
\end{tikzpicture}
\caption{Pierwsze płatne wdrożenie --- sprawdzian, czy wiedza nie zależy od jednej osoby.}
\end{figure}


\begin{enumerate}
  \item \textbf{Obsada.} Jeden Założyciel odpowiada za dostarczenie pilota i uzgadnia odbiór z klientem. Drugi prowadzi wydanie, a wyznaczony zastępca ćwiczy odtworzenie. Pozostali mogą dzielić pracę Kuźni, ocenę eksperymentów i koordynację testów. Pierwszą osobę z otwartego naboru przyjmuje się do Ekipy Produktowej i powierza jej podsystem, który będzie później utrzymywać.
  \item \textbf{Rozpoznanie.} Ekipa od Zadań Otwartych dostaje tydzień na sprawdzenie, czy obraz o zmierzchu wystarcza do osiągnięcia uzgodnionego celu. Sieć Terenowa wykonuje płatne próby w kilku lokalizacjach według tej samej karty.
  \item \textbf{Rozwiązanie.} Jeśli próby wykażą powtarzalny problem z kalibracją, Ekspert Basiliska otrzymuje mandat na opracowanie metody. Partner Specjalistyczny może wykonać przyrząd według specyfikacji. Kuźnia odbiera go na własnym stanowisku i przygotowuje procedurę dla Ekipy Produktowej.
  \item \textbf{Przekazanie.} Nowo zatrudniony inżynier samodzielnie instaluje wydanie na drugim urządzeniu, przeprowadza próbę odbiorową i zapisuje wynik. Opiekunowie Systemów sprawdzają powrót do poprzedniej wersji. Kierownik omawia z klientem wynik i wskazuje kontakt do utrzymania.
\end{enumerate}

\begin{Prawo}
\begin{itemize}
  \item \textbf{Umowa z klientem:} licencja według Ogólnych warunków licencyjnych i modułu właściwego dla produktu; Zamówienie podpisują osoby uprawnione do reprezentacji, własnoręcznie albo kwalifikowanym podpisem elektronicznym. Przed licencją --- lista kontrolna z Załącznika D (klasyfikacja eksportowa, sankcje, oświadczenie użytkownika końcowego).
  \item \textbf{Dane klienta:} jeżeli wdrożenie obejmuje dane osobowe --- umowa powierzenia (art.~28 RODO); Dane Operacyjne do rozwoju --- wyłącznie po anonimizacji.
  \item \textbf{Klient z sektora obronnego:} przed pierwszymi rozmowami --- wniosek o świadectwo bezpieczeństwa przemysłowego i ocena wymogu koncesji; skład zespołu wdrożeniowego z uwzględnieniem poświadczeń bezpieczeństwa.
  \item \textbf{Prawa do przyrządu Partnera:} odbiór przyrządu obejmuje dokumentację, pliki źródłowe i prawa do serwisowania (wzór~2) --- inaczej Spółka nie utrzyma rozwiązania u klienta.
\end{itemize}
\end{Prawo}

Po pilocie Założyciele przeglądają trzy dane: które czynności wciąż potrafi wykonać tylko jedna osoba, ile czasu zajęła druga instalacja i które problemy terenowe pozostają niewyjaśnione. Te odpowiedzi pomagają wybrać kolejny etat, grant albo zlecenie.

% ================================================================
\section{Strategia antymobbingowa}
\label{sec:mobbing}

Od 5 listopada 2026 r. obowiązuje nowe brzmienie art.~94\textsuperscript{3} Kodeksu pracy (Dz.U. z 2026 r. poz.~1046). Pracodawca ma nie tylko „przeciwdziałać mobbingowi”, ale robić to systematycznie, w pięciu obszarach:
\begin{quote}
„Pracodawca jest obowiązany systematycznie przeciwdziałać mobbingowi w szczególności przez działania prewencyjne, wykrywanie mobbingu oraz właściwe reagowanie na mobbing, a także przez działania naprawcze i wsparcie osób dotkniętych mobbingiem.” (art.~94\textsuperscript{3} §~1 KP)
\end{quote}
Strategia Basiliska jest zbudowana wokół tych pięciu obszarów. Obejmuje także szerszy obowiązek przeciwdziałania nierównemu traktowaniu, dyskryminacji i naruszaniu godności oraz innych dóbr osobistych pracownika (art.~94 pkt~2b--2c KP).

\subsection*{7.1. Co jest, a co nie jest mobbingiem}
\begin{itemize}
  \item \textbf{Mobbing} to uporczywe nękanie pracownika, czyli nękanie powtarzalne, nawracające albo stałe. Przejawami są w szczególności: upokarzanie, uwłaczanie, zastraszanie, zaniżanie oceny przydatności zawodowej, nieuzasadniona krytyka, poniżanie albo ośmieszanie, utrudnianie wykonywania zadań, komunikacji albo dostępu do informacji, izolowanie i eliminowanie z zespołu (art.~94\textsuperscript{3} §~2--4 KP).
  \item \textbf{Cel sprawcy nie ma znaczenia}: zachowania są mobbingiem, „choćby ich celem nie było uporczywe nękanie” (§~7). Mobbingiem jest też nakazanie takich zachowań albo zachęcanie do nich (§~6).
  \item \textbf{Sprawcą może być każdy} --- pracodawca, przełożony, współpracownik, podwładny, a także osoby pracujące na innej podstawie niż umowa o pracę, czyli konsultanci B2B, partnerzy i testerzy (§~8).
  \item \textbf{Nie jest mobbingiem} uzasadnione i wyrażone we właściwej formie rozliczanie z pracy albo jej krytyka (§~9) --- w szczególności rzeczowy przegląd kodu, ocena wyniku próby albo odmowa odbioru niespełniającego kryteriów. Nie jest nim też zachowanie incydentalne (§~3), które jednak może naruszać dobra osobiste i również wymaga reakcji (art.~94 pkt~2c KP).
\end{itemize}

\textbf{Skutki dla Spółki:} pracownik, który doznał mobbingu, może żądać zadośćuczynienia w wysokości \textbf{co najmniej sześciokrotności minimalnego wynagrodzenia} albo odszkodowania (§~11). Spółka może następnie dochodzić zwrotu od sprawcy, w stopniu odpowiadającym jego winie (§~12).

\begin{Prawo}
\textbf{Ryzyka specyficzne dla Basiliska.} W małym zespole Założyciele są jednocześnie przełożonymi, dyrektorami i większościowymi akcjonariuszami, co tworzy silną asymetrię władzy. Presję zwiększają terminy pilotów, dyżury Opiekunów Systemów, próby terenowe i kultura szczegółowego przeglądu pracy. Strategia przewiduje dlatego niezależną ścieżkę zgłoszeń wobec Założycieli i dyrektorów (pkt~7.3) oraz zasady przeglądu pracy, które odróżniają krytykę wyniku od oceny osoby (pkt~7.2).
\end{Prawo}

\subsection*{7.2. Prewencja}
\begin{enumerate}
  \item \textbf{Zasady współpracy} przyjęte przez Radę Dyrektorów i podpisywane przy rozpoczęciu współpracy przez każdą osobę, w tym Założycieli, konsultantów B2B, Partnerów i testerów. Stanowią część regulaminu (Załącznik).
  \item \textbf{Przegląd pracy dotyczy wyniku, nie osoby.} Uwagi w przeglądzie kodu, dokumentacji i prób odbiorczych formułuje się rzeczowo, z odwołaniem do kryteriów. Polecenia i oceny o istotnym znaczeniu dla pracownika przekazuje się w formie dokumentowej.
  \item \textbf{Obiektywne kryteria oceny i awansu} (art.~94 pkt~9 KP) --- jawne kryteria awansu, grantów eksperckich i premii z rozdziału~4 ograniczają ryzyko zarzutu „zaniżania oceny przydatności zawodowej”.
  \item \textbf{Organizacja pracy:} sprawiedliwa rotacja dyżurów, limit czasu pracy przy pilotach, prawo do odpoczynku (art.~132--133 KP), dostęp do informacji potrzebnych do pracy nadawany według zadań, a nie według relacji osobistych.
  \item \textbf{Szkolenie} przy rozpoczęciu współpracy i co najmniej raz w roku dla całego zespołu; dla Założycieli, liderów Ekip i Opiekunów Systemów --- szkolenie z prowadzenia trudnych rozmów i udzielania informacji zwrotnej.
  \item \textbf{Anonimowe badanie atmosfery} raz w roku i rozmowa przy zakończeniu współpracy; wyniki zbiorcze otrzymuje Rada Dyrektorów.
\end{enumerate}

\begin{DlaHR}
\textbf{Roczny kalendarz HR dla prewencji:} styczeń --- przegląd regulaminu i raport Pełnomocnika dla Rady; marzec --- szkolenie całego zespołu; czerwiec --- anonimowe badanie atmosfery (wyniki zbiorcze, bez grup mniejszych niż 5 osób); wrzesień --- szkolenie z informacji zwrotnej dla Założycieli i liderów; na bieżąco --- szkolenie przy rozpoczęciu współpracy, rozmowy przy odejściu i przegląd grafiku dyżurów pod kątem równego obciążenia.
\end{DlaHR}

\subsection*{7.3. Wykrywanie --- kanały zgłoszeń}
\begin{enumerate}
  \item \textbf{Pełnomocnik ds. przeciwdziałania mobbingowi} wyznaczony uchwałą Rady, z zastępcą. Nie powinien nim być dyrektor wykonawczy ani bezpośredni przełożony większości zespołu.
  \item \textbf{Dyrektor niewykonawczy} --- dla zgłoszeń dotyczących dyrektora wykonawczego albo Pełnomocnika.
  \item \textbf{Zewnętrzny doradca} (np. kancelaria albo psycholog organizacji) --- dla zgłoszeń dotyczących Założyciela albo dyrektora, a także gdy zgłaszający nie ufa kanałom wewnętrznym. Dopuszcza się zgłoszenie anonimowe.
  \item Zgłoszenie może dokonać osoba, której dotyczy nękanie, albo świadek, w dowolnej formie; przyjęcie potwierdza się w ciągu 3 dni roboczych.
\end{enumerate}

\subsection*{7.4. Reagowanie --- postępowanie wyjaśniające}
\begin{enumerate}
  \item \textbf{Ocena wstępna} w ciągu 7 dni: czy zgłoszenie dotyczy mobbingu, dyskryminacji, naruszenia dóbr osobistych albo innego problemu (konflikt, organizacja pracy).
  \item \textbf{Komisja} trzyosobowa: Pełnomocnik, osoba wskazana przez przedstawicieli pracowników i osoba niezależna. W składzie nie może być osoba, której dotyczy zgłoszenie, jej podwładny ani przełożony. Jeżeli zgłoszenie dotyczy Założyciela albo dyrektora, komisję powołują dyrektorzy niewykonawczy, a jej przewodniczącym jest zewnętrzny doradca; dyrektor, którego dotyczy sprawa, nie uczestniczy w żadnej decyzji (\US{\ref{art:rada} ust.~6}).
  \item \textbf{Środki tymczasowe:} rozdzielenie stron, zmiana podległości albo trybu pracy --- bez pogorszenia warunków osoby zgłaszającej.
  \item \textbf{Prawa stron:} wysłuchanie obu stron, możliwość udziału osoby zaufania, poufność, ochrona danych (art.~6 ust.~1 lit.~c i art.~9 ust.~2 lit.~b RODO).
  \item \textbf{Termin:} 30 dni od powołania komisji, z możliwością jednorazowego przedłużenia o 30 dni. Komisja sporządza protokół z ustaleniami i rekomendacją; decyzję o działaniach podejmuje Rada Dyrektorów (bez udziału dyrektora, którego dotyczy sprawa).
\end{enumerate}

\begin{figure}[H]
\centering
\begin{tikzpicture}[node distance=7mm and 8mm]
  \node[box=4.6cm, fill=black!8] (z) {\textbf{Zgłoszenie}\\Pełnomocnik, dyrektor niewykonawczy albo doradca zewnętrzny; także anonimowo\\potwierdzenie w 3 dni robocze};
  \node[q=3.6cm, below=of z] (q1) {Ocena wstępna (7 dni): mobbing, dyskryminacja, naruszenie dóbr?};
  \node[box=3.6cm, right=10mm of q1] (inne) {\textbf{Inne działania}\\rozmowa, mediacja, zmiana organizacji pracy};
  \node[q=3.6cm, below=of q1] (q2) {Dotyczy Założyciela albo dyrektora?};
  \node[xx=3.6cm, right=10mm of q2] (ext) {Komisję powołują dyrektorzy niewykonawczy; przewodniczy doradca zewnętrzny};
  \node[rd=4.6cm, below=of q2] (kom) {\textbf{Komisja (3 osoby)}\\wysłuchanie stron, osoba zaufania, poufność\\protokół w 30 dni (+30)};
  \node[dw=3.2cm, left=10mm of kom] (tymcz) {\textbf{Środki tymczasowe}\\rozdzielenie stron bez szkody dla zgłaszającego};
  \node[rd=4.6cm, below=of kom] (rd) {\textbf{Decyzja Rady Dyrektorów}\\bez udziału dyrektora, którego dotyczy sprawa};
  \node[xx=3.4cm, below=of rd] (s2) {\textbf{Wsparcie}\\psycholog, zmiana organizacji pracy};
  \node[xx=3.4cm, left=6mm of s2] (s1) {\textbf{Wobec sprawcy}\\kara, rozwiązanie umowy, regres};
  \node[wz=3.4cm, right=6mm of s2] (s3) {\textbf{Raport roczny}\\wnioski i zmiany systemowe};
  \draw[line] (z) -- (q1);
  \draw[line] (q1) -- node[lbl, above]{nie} (inne);
  \draw[line] (q1) -- node[lbl, left]{tak} (q2);
  \draw[line] (q2) -- node[lbl, above]{tak} (ext);
  \draw[line] (q2) -- node[lbl, left]{nie} (kom);
  \draw[line] (ext) |- (kom);
  \draw[line, dashed] (kom) -- (tymcz);
  \draw[line] (kom) -- (rd);
  \draw[line] (rd) -- (s2); \draw[line] (rd) -| (s1); \draw[line] (rd) -| (s3);
\end{tikzpicture}
\caption{Postępowanie po zgłoszeniu mobbingu (art.~94\textsuperscript{3} §~1 KP; regulamin, §§~6--8).}
\end{figure}

\begin{DlaHR}
\textbf{Jak przyjąć zgłoszenie:}
\begin{enumerate}
  \item Wysłuchaj bez oceniania i bez przerywania; zapytaj o fakty: co, kiedy, kto, kto był świadkiem, czy są wiadomości albo dokumenty.
  \item Zapisz relację możliwie dosłownie, z datą; daj zgłaszającemu do przeczytania.
  \item Nie obiecuj wyniku ani pełnej anonimowości wobec komisji --- wyjaśnij, kto pozna sprawę i jakie są terminy (potwierdzenie 3 dni robocze, ocena wstępna 7 dni, postępowanie 30 dni).
  \item Nie konfrontuj samodzielnie stron i nie informuj osoby wskazanej przed decyzją Pełnomocnika.
  \item Jeżeli zgłaszający jest w złym stanie albo grozi mu dalsze nękanie --- od razu zaproponuj środki tymczasowe i kontakt z psychologiem.
  \item Przekaż zgłoszenie Pełnomocnikowi najpóźniej następnego dnia roboczego; gdy dotyczy Założyciela albo dyrektora --- dyrektorowi niewykonawczemu albo doradcy zewnętrznemu.
  \item Dokumentację przechowuj oddzielnie od akt osobowych, z dostępem tylko dla Pełnomocnika i komisji.
\end{enumerate}
Pamiętaj: rzeczowa krytyka pracy nie jest mobbingiem, ale pojedyncze niewłaściwe zachowanie też wymaga reakcji --- jako naruszenie dóbr osobistych.
\end{DlaHR}

\subsection*{7.5. Działania naprawcze i wsparcie}
\begin{itemize}
  \item \textbf{Wobec sprawcy-pracownika:} rozmowa dyscyplinująca, kara porządkowa upomnienia albo nagany (art.~108 KP), wypowiedzenie albo rozwiązanie umowy bez wypowiedzenia (art.~52 KP); roszczenie zwrotne Spółki (art.~94\textsuperscript{3} §~12 KP).
  \item \textbf{Wobec konsultanta, Partnera albo testera:} rozwiązanie umowy z ważnych powodów (art.~746 KC), na podstawie klauzuli w umowie.
  \item \textbf{Wobec Założyciela albo dyrektora:} odwołanie z Rady Dyrektorów przez Walne Zgromadzenie (bezwzględna większość głosów oddanych), rozwiązanie umowy, roszczenie zwrotne. Uwaga na skutek dla vestingu --- patrz ramka poniżej.
  \item \textbf{Wsparcie osoby dotkniętej:} konsultacje psychologiczne finansowane przez Spółkę, zmiana organizacji pracy na jej wniosek, monitorowanie sytuacji przez 6 miesięcy. Zakaz działań odwetowych wobec zgłaszającego i świadków; zgłoszenie w dobrej wierze nie może być podstawą negatywnej oceny ani rozwiązania umowy.
  \item \textbf{Wnioski systemowe:} raz w roku Pełnomocnik przedstawia Radzie Dyrektorów zbiorczy raport (liczba zgłoszeń, czas postępowań, działania, wyniki badania atmosfery) i propozycje zmian.
\end{itemize}

\begin{Prawo}
\textbf{Luka w umowie spółki --- mobbing przez Założyciela.} Zamknięta lista „istotnych przyczyn” z \US{\ref{art:obowiazek-zbycia} ust.~5} obejmuje wyłącznie: prawomocne skazanie za umyślne przestępstwo przeciwko Spółce, naruszenie obowiązków dotyczących IP, poufności, zakazu konkurencji albo kontroli eksportu oraz uporczywe niewykonywanie obowiązków. \textbf{Mobbingu na niej nie ma.} Jeżeli Spółka odwoła Założyciela z powodu mobbingu, będzie to Odejście Uzasadnione --- z dodatkowym zwolnieniem akcji za 12 miesięcy --- i bez zakazu konkurencji z \UA{\ref{sha:konkurencja} ust.~1}.

Zamknięcie tej luki \textbf{wymaga zmiany umowy spółki}: dodania w \US{\ref{art:obowiazek-zbycia} ust.~5} punktu „d)~mobbing albo dyskryminacja stwierdzone w postępowaniu wyjaśniającym prowadzonym według regulaminu Spółki z udziałem niezależnego przewodniczącego albo prawomocnym wyrokiem”. Ponieważ zmiana zaostrza obowiązki związane z akcjami, wymaga --- poza 75\% wszystkich głosów i protokołem notarialnym --- zgody każdego Założyciela, którego dotyczy (\US{\ref{art:final} ust.~1}). Zalecenie: uzgodnić ją teraz, przed aktem albo przed pierwszą rundą, gdy żadna sprawa nie jest w toku. Do tego czasu strategia działa bez tej zmiany, a Spółka zachowuje roszczenie zwrotne i możliwość odwołania dyrektora.
\end{Prawo}

\subsection*{7.6. Regulamin --- kiedy i jak}
\begin{itemize}
  \item \textbf{Obowiązek przeciwdziałania mobbingowi} dotyczy Spółki od pierwszego pracownika (art.~94\textsuperscript{3} §~1 KP).
  \item \textbf{Regulamin} określający reguły, procedury i częstotliwość działań jest obowiązkowy od \textbf{10 pracowników} (art.~94\textsuperscript{3a} §~1 KP). Spółka przyjmuje go dobrowolnie wcześniej --- według wzoru z Załącznika --- uchwałą Rady Dyrektorów.
  \item Po osiągnięciu progu 10 pracowników treść regulaminu uzgadnia się z przedstawicielami pracowników wyłonionymi w trybie przyjętym w Spółce, jeżeli nie działa związek zawodowy; po 30 dniach bez uzgodnienia Spółka ustala treść sama, uwzględniając dotychczasowe ustalenia (art.~94\textsuperscript{3a} §~4--5 KP). Regulamin wchodzi w życie po upływie 2 tygodni od podania go do wiadomości pracowników (art.~104\textsuperscript{3} KP w zw. z art.~94\textsuperscript{3a} §~6 KP).
  \item Regulamin obejmuje także osoby pracujące na innej podstawie niż umowa o pracę --- przez klauzulę w ich umowach.
\end{itemize}

% ================================================================
\section{Czy trzeba zmieniać umowę spółki?}
\label{sec:zmiany}

\textbf{Nie.} Umowa spółki określa organy, ich kompetencje, progi decyzyjne, akcje i prawa własności intelektualnej, ale nie reguluje struktury zespołu, stanowisk ani form zatrudnienia. Siedem ról z tej strategii to organizacja pracy w ramach kompetencji Rady Dyrektorów, która prowadzi sprawy Spółki (\US{\ref{art:rada} ust.~2}) i przyjmuje budżet (\US{\ref{art:rada} ust.~7}).

\begin{center}
\small
\begin{tabularx}{\linewidth}{@{}>{\raggedright\arraybackslash}X >{\raggedright\arraybackslash}p{3.6cm} >{\raggedright\arraybackslash}p{4.0cm}@{}}
\toprule
\textbf{Działanie} & \textbf{Dokument / organ} & \textbf{Zmiana umowy spółki?} \\
\midrule
Utworzenie ról, zasady awansu, premii i grantów eksperckich & uchwała Rady; regulamin Rady (\US{\ref{art:rada} ust.~8}) & nie \\
Budżet zatrudnienia, testów i zleceń & uchwała Rady (budżet roczny) & nie \\
Delegacja zatrudniania dla dyrektora wykonawczego & uchwała Rady (\US{\ref{art:rada} ust.~4}) & nie \\
Pełnomocnictwo dla kierownika wdrożenia do podpisywania zamówień i protokołów do określonej kwoty & pełnomocnictwo udzielone przez dwóch dyrektorów (art.~98 KC) & nie \\
Prokura dla menedżera & zgoda wszystkich dyrektorów (\US{\ref{art:reprezentacja} ust.~2}) & nie \\
Opcje dla zespołu (seria P) w trybie zwykłej emisji & uchwała WZ 75\% (\US{\ref{art:seria-p} ust.~2}) & nie \\
Opcje w trybie warunkowej emisji (regulamin programu) & uchwała WZ 4/5 i 75\% z protokołem notarialnym & formalnie tak --- uchwała o warunkowej emisji zmienia umowę spółki, ale umowa już ją przewiduje (\US{\ref{art:seria-p} ust.~2}) \\
Pula opcji większa niż 176\,471 akcji & uchwała WZ 75\% & \textbf{tak} (\US{\ref{art:seria-p} ust.~3}) \\
Nowa osoba o roli Założyciela (seria F) & dwie uchwały WZ 75\% (\US{\ref{art:seria-f}}) & nie --- w limicie 250\,000 akcji \\
Role Założycieli w nowej strukturze & aktualizacja tabeli w \UA{\ref{sha:zaangazowanie} ust.~1} & nie (zmiana UA, zgoda Stron) \\
Wynagrodzenie dyrektorów & uchwała WZ (\US{\ref{art:wz-zastrzezone} ust.~3}) & nie \\
Monitoring dzienników dostępu & polityka bezpieczeństwa + obwieszczenie (art.~22\textsuperscript{3} KP) & nie \\
Strategia i regulamin antymobbingowy, Pełnomocnik, komisja & uchwała Rady; od 10 pracowników --- uzgodnienie z przedstawicielami pracowników (art.~94\textsuperscript{3a} KP) & nie \\
Klauzula antymobbingowa w umowach B2B, z Partnerami i testerami & umowy (wzór~B, wzór zlecenia) & nie \\
Mobbing Założyciela jako „istotna przyczyna” odejścia (vesting) & uchwała WZ 75\%, protokół notarialny, zgoda każdego Założyciela & \textbf{tak} (\US{\ref{art:obowiazek-zbycia} ust.~5}) --- zalecane, nieobowiązkowe \\
\bottomrule
\end{tabularx}
\end{center}

\textbf{Do przygotowania przed wdrożeniem strategii} (bez zmiany umowy spółki):
\begin{enumerate}
  \item uchwała Rady przyjmująca strategię, role, zasady awansu, premii i grantów oraz budżet zatrudnienia;
  \item regulamin Rady Dyrektorów z delegacją zatrudniania dla dyrektora wykonawczego i zasadami udzielania pełnomocnictw;
  \item wzór umowy zlecenia dla Sieci Terenowej z przeniesieniem praw do danych i klauzulą RODO;
  \item krótka umowa na płatne zadanie rekrutacyjne;
  \item obwieszczenie o monitoringu (przed zatrudnieniem pierwszego pracownika);
  \item uchwały Walnego Zgromadzenia w sprawie programu motywacyjnego, jeżeli opcje mają być przyznawane od początku;
  \item opcjonalnie: aktualizacja tabeli ról Założycieli w \UA{\ref{sha:zaangazowanie} ust.~1};
  \item uchwała Rady przyjmująca regulamin antymobbingowy (Załącznik), wyznaczenie Pełnomocnika i wybór zewnętrznego doradcy;
  \item klauzula antymobbingowa w umowach z konsultantami, Partnerami i testerami.
\end{enumerate}

\textbf{Jedyny punkt, który wymagałby zmiany umowy spółki,} to dopisanie mobbingu i dyskryminacji do zamkniętej listy istotnych przyczyn odejścia Założyciela (rozdział~\ref{sec:mobbing}, pkt~7.5). Strategia działa bez tej zmiany; zmiana jest zalecana, ale wymaga zgody wszystkich Założycieli.

% ================================================================
\section{Wzory do skopiowania}

\subsection*{Wzór 1. Karta mandatu eksperta lub zespołu rozpoznawczego}
\begin{tabularx}{\linewidth}{@{}>{\raggedright\arraybackslash}p{3.4cm} X@{}}
\toprule
\textbf{Pole} & \textbf{Do uzupełnienia} \\
\midrule
Typ i pytanie & grant ekspercki / krótkie rozpoznanie; jaki problem rozwiązujemy lub jaką hipotezę sprawdzamy \\
Obsada & osoba prowadząca, sponsor decyzji, osoba przejmująca wynik \\
Podstawa & aneks do umowy o pracę nr \field{\ldots} / zlecenie nr \field{\ldots}; Umowa nr~2 (osoby spoza zatrudnienia) \\
Zasoby & okres, chroniony czas, sprzęt, limit wydatków (w budżecie --- dyrektor wykonawczy; poza budżetem ponad 250\,000 zł --- Rada), zakres samodzielnych decyzji \\
Informacje & kategoria dostępu według polityki; informacje Kontrolowane --- zgoda Rady z dnia \field{\ldots} \\
Dowód & dla eksperta --- wynik do odtworzenia; dla rozpoznania --- obserwacja uzasadniająca kontynuację lub zamknięcie \\
Rozliczenie & wynagrodzenie, premia za rezultat (kryterium i termin), termin przeglądu \\
Zamknięcie & miejsce zapisania materiałów w systemach Spółki, odbiorca, decyzja o przekazaniu, odnowieniu lub zakończeniu, odebranie dostępów \\
\bottomrule
\end{tabularx}

\subsection*{Wzór 2. Karta zlecenia dla partnera specjalistycznego}
\begin{tabularx}{\linewidth}{@{}>{\raggedright\arraybackslash}p{3.4cm} X@{}}
\toprule
\textbf{Pole} & \textbf{Do uzupełnienia} \\
\midrule
Strony i cel & wykonawca, inżynier odbierający po stronie Basiliska, rezultat i zakres; podstawa: umowa nr \field{\ldots} (wzór~B / umowa o dzieło) \\
Zatwierdzenie & w budżecie / uchwała Rady nr \field{\ldots}; podmiot powiązany: tak / nie \\
Dostawa & terminy i etapy, dokumentacja, pliki źródłowe, instrukcja serwisowa, wykaz komponentów (licencja, kraj pochodzenia, klasyfikacja eksportowa) \\
Odbiór & próba odbiorowa, tolerancje, termin oceny i procedura poprawek; protokół odbioru podpisują osoby uprawnione do reprezentacji albo pełnomocnik \\
Rozliczenie & cena, wydatki, terminy płatności, rozliczenie w razie wcześniejszego zakończenia \\
Dalsze użycie & przeniesienie praw (Umowa nr~2), licencja na elementy wcześniejsze z prawem do sublicencji, przekazanie wiedzy, zamknięcie dostępów \\
\bottomrule
\end{tabularx}

\subsection*{Wzór 3. Karta płatnego testu terenowego}
\begin{tabularx}{\linewidth}{@{}>{\raggedright\arraybackslash}p{3.4cm} X@{}}
\toprule
\textbf{Pole} & \textbf{Do uzupełnienia} \\
\midrule
Podstawa & umowa ramowa zlecenia nr \field{\ldots} z dnia \field{\ldots} \\
Warunki udziału & urządzenie, miejsce lub typ lokalizacji, termin, wymagane zgody (właściciel terenu, przestrzeń powietrzna, kategoria lotu) \\
Instrukcja & kroki próby, warunek przerwania, kontakt do koordynatora \\
Dane & co zebrać, czego unikać (wizerunek osób trzecich, obiekty chronione), gdzie przesłać plik, kto może go wykorzystać i jak długo będzie przechowywany \\
Kryterium jakości & wymagane informacje o warunkach próby, format, sposób udokumentowania usterki \\
Zapłata & stawka (nie niższa niż minimalna stawka godzinowa), uznawane wydatki, zasady odwołania zadania i termin płatności \\
Weryfikacja & osoba sprawdzająca, termin odpowiedzi, możliwość poprawienia niekompletnego zgłoszenia \\
\bottomrule
\end{tabularx}

% ================================================================
\section*{Załącznik --- Wzór regulaminu przeciwdziałania mobbingowi, dyskryminacji i naruszaniu dóbr osobistych}
\addcontentsline{toc}{section}{Załącznik --- Wzór regulaminu antymobbingowego}

{\small Przyjmuje Rada Dyrektorów uchwałą nr \field{numer} z dnia \field{data}. Od 10 pracowników --- po uzgodnieniu z przedstawicielami pracowników (art.~94\textsuperscript{3a} KP).\par}

\begin{enumerate}
  \item[§~1.] \textbf{Zakres.} Regulamin określa reguły, procedury i częstotliwość działań Basilisk Systems P.S.A. w zakresie przeciwdziałania mobbingowi, dyskryminacji, nierównemu traktowaniu i naruszaniu godności oraz innych dóbr osobistych. Stosuje się go do pracowników oraz --- na podstawie postanowień umów --- do dyrektorów, konsultantów, Partnerów, testerów i innych osób wykonujących pracę na rzecz Spółki.
  \item[§~2.] \textbf{Definicje.} Mobbing, dyskryminację i molestowanie rozumie się zgodnie z art.~94\textsuperscript{3} i 18\textsuperscript{3a} KP. Uzasadniona i wyrażona we właściwej formie krytyka pracy nie jest mobbingiem.
  \item[§~3.] \textbf{Zasady współpracy.} Każda osoba objęta regulaminem: traktuje innych z szacunkiem; formułuje uwagi rzeczowo i wobec wyniku pracy, nie osoby; nie wyklucza nikogo z komunikacji ani z dostępu do informacji potrzebnych do pracy; reaguje albo zgłasza zachowania naruszające regulamin.
  \item[§~4.] \textbf{Pełnomocnik.} Rada Dyrektorów wyznacza Pełnomocnika ds. przeciwdziałania mobbingowi i jego zastępcę oraz wskazuje zewnętrznego doradcę. Dane kontaktowe są dostępne dla wszystkich.
  \item[§~5.] \textbf{Prewencja i częstotliwość.} Szkolenie przy rozpoczęciu współpracy i co najmniej raz w roku; anonimowe badanie atmosfery raz w roku; rozmowa przy zakończeniu współpracy; roczny raport Pełnomocnika dla Rady Dyrektorów; przegląd regulaminu co najmniej raz w roku.
  \item[§~6.] \textbf{Zgłoszenia.} Zgłoszenie może złożyć każda osoba --- ustnie, pisemnie albo elektronicznie --- Pełnomocnikowi, dyrektorowi niewykonawczemu albo zewnętrznemu doradcy, także anonimowo. Przyjęcie potwierdza się w ciągu 3 dni roboczych.
  \item[§~7.] \textbf{Postępowanie.} Ocena wstępna w ciągu 7 dni; komisja trzyosobowa bez udziału osób powiązanych ze sprawą; w sprawach dotyczących Założyciela albo dyrektora --- komisja powołana przez dyrektorów niewykonawczych pod przewodnictwem zewnętrznego doradcy. Wysłuchanie stron, osoba zaufania, poufność. Protokół i rekomendacja w ciągu 30 dni (przedłużenie o 30 dni). Decyzję podejmuje Rada Dyrektorów bez udziału dyrektora, którego dotyczy sprawa.
  \item[§~8.] \textbf{Środki.} Środki tymczasowe bez pogorszenia sytuacji zgłaszającego; środki wobec sprawcy według Kodeksu pracy albo umowy; wsparcie psychologiczne i organizacyjne dla osoby dotkniętej; roszczenie zwrotne Spółki.
  \item[§~9.] \textbf{Ochrona przed odwetem.} Zgłoszenie dokonane w dobrej wierze ani udział w postępowaniu nie mogą być przyczyną niekorzystnego traktowania. Zgłoszenie świadomie fałszywe stanowi naruszenie obowiązków.
  \item[§~10.] \textbf{Dokumentacja i dane.} Dokumentację postępowań przechowuje Pełnomocnik oddzielnie od akt osobowych, z ograniczonym dostępem, przez okres przedawnienia roszczeń. Dane przetwarza się w zakresie niezbędnym do wykonania obowiązku z art.~94\textsuperscript{3} KP.
  \item[§~11.] \textbf{Wejście w życie.} Regulamin wchodzi w życie po upływie 2 tygodni od podania go do wiadomości osób nim objętych.
\end{enumerate}

\vspace{1em}
{\small Dokument nie zmienia umowy spółki ani umowy akcjonariuszy; w razie rozbieżności rozstrzygają te umowy. Kwestie oznaczone jako wymagające weryfikacji (minimalna stawka godzinowa, składki testerów, zgody dotyczące lotów) należy potwierdzić z doradcą przed pierwszym naborem.\par}

\end{document}
